\documentclass{article}%
\usepackage[T1]{fontenc}%
\usepackage[utf8]{inputenc}%
\usepackage{lmodern}%
\usepackage{textcomp}%
\usepackage{lastpage}%
\usepackage{geometry}%
\geometry{tmargin=5cm,lmargin=5cm}%
\usepackage{graphicx}%
\usepackage{float}
%
%
%
\begin{document}%
\normalsize%
\section{Model Outputs:}%
\label{sec:PlotofFlattops}%
\begin{figure}[H]%
\centering%
\includegraphics[width=\textwidth]{/home/maniacalm/bin/VQVAE_studies/run_data/flattops_init/images/outputs.png}%
\caption{Outputs for the models for certain SNR datasets}%
\end{figure}
%
Models were trained with synthetic datasets with an SNR around the specified value but with variation of a Gaussian with $\sigma = \frac{snr}{5*2}$. So $95\%$ of the runs in each dataset are within $20\%$ of the snr. In hindsight this is far too precise of datasets for this but for the goal of looking at the performance of n\_codes on model behavior is more of a side note.

\begin{figure}[H]%
\centering%
\includegraphics[width=\textwidth]{/home/maniacalm/bin/VQVAE_studies/run_data/flattops_init/images/n_codes_vs_flattness_scatter.png}%
\caption{Comparison N\_codes and the flattness metric with size denoting the relative loss}%
\end{figure}
%
To measure the flatness seen in other runs a metric was used that should manage to track relative flatness of the outputs \[flattness = \int (f(x))''dx\] with $f(x)$ being the data. The next two plots being the same but for the comparing relative to the loss instead.

\begin{figure}[H]%
\centering%
\includegraphics[width=\textwidth]{/home/maniacalm/bin/VQVAE_studies/run_data/flattops_init/images/n_codes_vs_loss.png}%
\end{figure}
%
\begin{figure}[H]%
\centering%
\includegraphics[width=\textwidth]{/home/maniacalm/bin/VQVAE_studies/run_data/flattops_init/images/n_codes_vs_flatness.png}%
\end{figure}

While no apparent overall pattern exists across the datasets. For the $snr=40$ dataset we see both flattness and loss drop with n\_codes. For the lower snr runs we see either a decrease or an unchanging effect. The snr of 5 runs had no clear pattern and should be looked at again.
%
\end{document}